\documentclass[main.tex]{subfiles}


\begin{document}

%from pagenumber to up
% \phantomsection 
% \hypertarget{toc}{}

 

 

% номер секции вручную
\setcounter{section}{3
}
\addtocounter{section}{-1} 

\section{Электризация}


\textbf{Электризация}~--- это процесс, в результате которого тело (часть тела) переходит из незаряженного состояния в заряженное. Можно выделить два вида электризации.
\begin{enumerate}
    \item При электризации \textbf{трением} тела приобретают заряды вследствие их трения друг о друга. 

    Пусть имеется два тела~--- незаряженные куски пластика~П и шерсти~Ш (рис.\,\ref{fig:p91}, слева).

    \begin{figure}[H]
    \centering

    \begin{tikzpicture}[font=\small,            line cap=round,                             >={Stealth[inset=0pt,round,length=3mm, angle'=20]},vec/.style={>={Stealth[inset=0pt,round,length=3.5mm, angle'=20]}}]


\tikzset{atom/.pic={
  \begin{scope}[shift={({rnd*360}:.05)}]
  \shade[ball color=red] (0,0) circle (\dn);
  \shade[ball color=NavyBlue!70,rotate={rnd*360}] ([shift={(180:{\dn cm})}]0,0) circle (\de);
  \end{scope}}
}

\tikzset{atompo/.pic={
  \begin{scope}[shift={({rnd*360}:.05)}]
  \shade[ball color=red] (0,0) circle (\dn);
  \end{scope}}
}

\tikzset{atomne/.pic={
  \begin{scope}[shift={({rnd*360}:.05)},rotate={rnd*360}]
  \shade[ball color=red] (0,0) circle (\dn);
  \shade[ball color=NavyBlue!70] ([shift={(180:{\dn cm})}]0,0) circle (\de);
  \shade[ball color=NavyBlue!70] ([shift={(0:{\dn cm})}]0,0) circle (\de);
  \end{scope}}
}

\def\dn{.15}
\def\de{.1}


%%%%%%%%%%%%%%%%%%%
%left pic
%%%%%%%%%%%%%%%%%%%

\begin{scope}[yshift=.5cm]

%q
% \node[above] at (1,2) {$q_\text{П}=0$};
% \node[above] at (4,2) {$q_\text{Ш}=0$};


%name
\node[above] at (1,2) {П};
\node[above] at (4,2) {Ш};


%solid
\fill[yellow, nearly transparent] (0,0) rectangle++ (2,2);
\draw[] (0,0) rectangle++ (2,2);

\fill[yellow, nearly transparent] (3,0) rectangle++ (2,2);
\draw[] (3,0) rectangle++ (2,2);

    
%atoms left
\foreach \y in {0,...,2}
\foreach \x in {0,...,2}
\pic[shift={(.4,.4)}] at ({0+(2-.8)/2*\x},{0+(2-.8)/2*\y}) {atom};


%%%

%atoms right
\foreach \y in {0,...,2}
\foreach \x in {0,...,2}
\pic[shift={(3.4,.4)}] at ({0+(2-.8)/2*\x},{0+(2-.8)/2*\y}) {atom};

\end{scope}


%%%%%%%%%%%%%%%%%%%
%right pic
%%%%%%%%%%%%%%%%%%%

\begin{scope}[xshift=8cm]

%name
\node[above] at (1,2) {П};
\node[above right] at (3,2.5) {Ш};


%solid
\fill[NavyBlue,nearly transparent] (0,0) rectangle++ (2,2);
\draw[] (0,0) rectangle++ (2,2);

\fill[red,nearly transparent] (2,.5) rectangle++ (2,2);
\draw[] (2,.5) --++ (2,0) --++ (0,2) --++ (-2,0) --++ (0,-.5);


%atoms left
\foreach \y in {0,...,2}
\foreach \x in {0,...,1}
\pic[shift={(.4,.4)}] at ({0+(2-.8)/2*\x},{0+(2-.8)/2*\y}) {atom};

\pic[shift={(.4,.4)}] at ({0+(2-.8)/2*2},{0+(2-.8)/2*0}) {atom};
\pic[shift={(.4,.4)}] at ({0+(2-.8)/2*2},{0+(2-.8)/2*1}) {atomne};
\pic[shift={(.4,.4)}] at ({0+(2-.8)/2*2},{0+(2-.8)/2*2}) {atomne};


%%%

%atoms right
\foreach \y in {0,...,2}
\foreach \x in {1,...,2}
\pic[shift={(2.4,.9)}] at ({0+(2-.8)/2*\x},{0+(2-.8)/2*\y}) {atom};

\pic[shift={(2.4,.9)}] at ({0+(2-.8)/2*0},{0+(2-.8)/2*2}) {atom};
\pic[shift={(2.4,.9)}] at ({0+(2-.8)/2*0},{0+(2-.8)/2*1}) {atompo};
\pic[shift={(2.4,.9)}] at ({0+(2-.8)/2*0},{0+(2-.8)/2*0}) {atompo};


%vec
\draw[->] (3,2.6) --++ (90:1);
\draw[->] (3,.4) --++ (-90:1);


\end{scope}


%\Longrightarrow
\path (6.5,1.5) node {$\Longrightarrow$};












    \end{tikzpicture}
\caption{Электризация трением}
\label{fig:p91}
\end{figure}
\nointerlineskip

    Тела~П и~Ш трут друг о друга (рис.\,\ref{fig:p91}, справа). В этом случае некоторые \emph{электроны} (синие шары), находящиеся в контактирующих слоях атомов и сравнительно слабо связанные с ядрами (красные шары) <<своего>> тела, переходят к ядрам другого тела~--- электроны стремятся к тем ядрам, с которыми они сильнее связаны. После разведения тела оказываются заряженными: тело~П, получившее электроны от тела~Ш, несет отрицательный заряд; тело~Ш, соответственно, имеет положительный заряд.

    \item При электризации \textbf{влиянием} части тела заряжаются вследствие перераспределения заряженных частиц внутри тела под действием другого тела.

    Пусть имеется одно незаряженное металлическое тело~М (рис.\,\ref{fig:p92}, слева).

        \begin{figure}[H]
    \centering

    \begin{tikzpicture}[font=\small,            line cap=round,                             >={Stealth[inset=0pt,round,length=3mm, angle'=20]},vec/.style={>={Stealth[inset=0pt,round,length=3.5mm, angle'=20]}}]


\tikzset{atom/.pic={
  \begin{scope}[shift={({rnd*360}:.025)}]
  \shade[ball color=red] (0,0) circle (\dn);
  \shade[ball color=NavyBlue!70,rotate={rnd*360}] ([shift={(180:{\dn*1cm +\de*1cm})}]0,0) circle (\de);
  \end{scope}}
}

\tikzset{atompo/.pic={
  \begin{scope}[shift={({rnd*360}:.025)}]
  \shade[ball color=red] (0,0) circle (\dn);
  \end{scope}}
}

\tikzset{atomne/.pic={
  \begin{scope}[shift={({rnd*360}:.025)},rotate={rnd*0}]
  \shade[ball color=red] (0,0) circle (\dn);
  \shade[ball color=NavyBlue!70] ([shift={(180+rand*90:{\dn*1cm +\de*1cm})}]0,0) circle (\de);
  \shade[ball color=NavyBlue!70] ([shift={(rand*90:{\dn*1cm +\de*1cm})}]0,0) circle (\de);
  \end{scope}}
}

\def\dn{.15}
\def\de{.1}


%%%%%%%%%%%%%%%%%%%
%left pic
%%%%%%%%%%%%%%%%%%%

\begin{scope}[yshift=.5cm]


%name
\node[above] at (1,2) {М};


%solid
\fill[lightgray, semitransparent] (0,0) rectangle++ (2,2);
\draw[] (0,0) rectangle++ (2,2);

    
%atoms left
\foreach \y in {0,...,2}
\foreach \x in {0,...,2}
\pic[shift={(.4,.4)}] at ({0+(2-.8)/2*\x},{0+(2-.8)/2*\y}) {atom};


\end{scope}


%%%%%%%%%%%%%%%%%%%
%right pic
%%%%%%%%%%%%%%%%%%%

\begin{scope}[xshift=5cm]

%name
\node[above] at (1.5,2) {З};
\node[above] at (4,2.5) {М};


\fill[red,nearly transparent] (3,.5) rectangle++ (1,2);
\fill[NavyBlue,nearly transparent] (4,.5) rectangle++ (1,2);
\draw[] (3,.5) --++ (2,0) --++ (0,2) --++ (-2,0) --++ (0,-2);


%circ
\shade[ball color=NavyBlue] (1.5,1.5) circle (.5);


%%%


%atoms right
\foreach \y in {0,...,2}
\foreach \x in {1,...,1}
\pic[shift={(3.4,.9)}] at ({0+(2-.8)/2*\x},{0+(2-.8)/2*\y}) {atom};

\pic[shift={(3.4,.9)}] at ({0+(2-.8)/2*0},{0+(2-.8)/2*2}) {atompo};
\pic[shift={(3.4,.9)}] at ({0+(2-.8)/2*0},{0+(2-.8)/2*1}) {atompo};
\pic[shift={(3.4,.9)}] at ({0+(2-.8)/2*0},{0+(2-.8)/2*0}) {atompo};

\pic[shift={(3.4,.9)}] at ({0+(2-.8)/2*2},{0+(2-.8)/2*0}) {atomne};
\pic[shift={(3.4,.9)}] at ({0+(2-.8)/2*2},{0+(2-.8)/2*1}) {atomne};
\pic[shift={(3.4,.9)}] at ({0+(2-.8)/2*2},{0+(2-.8)/2*2}) {atomne};




\end{scope}


%\Longrightarrow
\path (4,1.5) node {$\Longrightarrow$};


    \end{tikzpicture}
\caption{Электризация влиянием}
\label{fig:p92}
\end{figure}
\nointerlineskip

    К телу~М подносят слева (не касаясь!), к примеру, отрицательно заряженный шар~З и устанавливают его некотором расстоянии от тела~М (рис.\,\ref{fig:p92}, справа). Шар~З отталкивает от себя \emph{свободные} электроны (синие шары) тела~М (часть электронов в металле являются свободными~--- то есть передвигающимися по всему объему тела подобно частицам газа в сосуде). 
    В результате в левой части тела~М протоны преобладают над электронами~--- левая часть тела заряжена положительно. Соответственно, правая часть этого тела от избытка электронов заряжена отрицательно. 

    % В результате часть электронов скапливается на правой поверхности тела~М, на левой поверхности остаются нескомпенсированные положительные заряды. При этом можно сказать, что в левой части тела~М протоны преобладают над электронами~--- левая часть тела заряжена положительно. Соответственно, правая часть тела от избытка электронов заряжена отрицательно 
    
\end{enumerate}



























 
\end{document}